% page 144 of the facsimile (image p-143 of the scan)
% block 157.5 x 237.7 mm ; left margin 8.0 mm
\pgc{-12.700mm}{0.000mm}{
\l{\cel{3}134.}
\l{}
\l{}
\l{\sou{emula}. Qua penas por egalesar o superesar la ceteri en kozi}
\l{\cel{3}laudinda od en konkurso loyala. - EFIS.}
\l{}
\l{\sou{emulsar}. (\sou{trans.}) Adjuntar a drinkajo liquido laktatra}
\l{\cel{3}quan onu extraktas de la semini ek qui onu povas ekpresar}
\l{\cel{3}oleo. - DEFIRS.}
\l{}
\l{\sou{emulsino}.\cel{2}Substanco albuminoida quan onu extraktas de la}
\l{\cel{3}mandeli.}
\l{}
\l{\sou{emundar}. (\sou{trans.})\cel{2}I. Purigar (arboro) de lua branchi mort-}
\l{\cel{3}inta, do lua planti parazita. - II. Netigar, beligar o}
\l{\cel{3}preparar kozo per depreno di omna parti o peci neutila,}
\l{\cel{3}nebona o nebela. - FIS.}
\l{}
\l{\sou{emuntorio}. (\sou{fiziol.}) Tubo, organo, per qua evakuesas la hu-}
\l{\cel{3}mori superflua.}
\l{}
\l{en. Qua esas inter limiti, spacale o tempale.}
\l{}
\l{\sou{enartroso}. (\sou{anat.}) Artiko movebla. - DEF.}
\l{}
\l{\sou{encefalo}. (\sou{anat.}) La ensemblo ek cerebro e cerebelo. - EFIS.}
\l{}
\l{\sou{encefalito}. (\sou{patol.}) Influro di la encefalo. - DEF.}
\l{}
\l{\sou{encikliko}. (\sou{eklezio katolika)}.\cel{2}Letro cirkulera da la papo,}
\l{\cel{3}pri punto di doktrino, di dogmato. - DEF.}
\l{}
\l{\sou{enciklopedio}. Libro qua kontenas expozita, segun ordino me-}
\l{\cel{3}todala, La toto di la savaji homala. - DEFIRS.}
\l{}
\l{\sou{-end-}. Sufixo qua signifikas \textquotedbl{}quan onu devas\textquotedbl{}, \textquotedbl{}quan oportas\textquotedbl{}.}
\l{}
\l{\sou{endemio}. (\sou{patol}.)morbo, di qua ula kazi renkontresas intermite}
\l{\cel{3}en ca o ta regiono. - DEFIRS.}
\l{}
\l{\sou{endemika}. Dicesas pri morbo qua esas specala ye regiono, ube}
\l{\cel{3}lu regnas preske sencese.- DEFIRS.}
\l{}
\l{\sou{endermika}. I. (\sou{medic.}) Di qua la fonto esas en la dermo.-}
\l{\cel{3}II. (\sou{pri metodo}). Segun qua onu efikas a la dermo per tri-}
\l{\cel{3}cioni, locioni, uncioni, e c.}
\l{}
\l{\sou{endetala}. (\sou{pri vendo o kompro di vari)}. Mikra-quantope. -}
\l{\cel{3}FDSued.}
\l{}
\l{\sou{endivio}. (\sou{bot.}) Cikorio gardenala. - L.\cel{1}\sou{cichorium endivia}. -F.}
\l{}
\l{\sou{endodermo}. (\sou{embriol.)}\cel{3}Strato interna di la blastodermo qua,}
\l{\cel{3}pose, divenos la digesto-tubo, e la kovro-epitelio e la}
\l{\cel{3}glandi anexa di ica.}
}
